\documentclass[10pt,a4paper]{amsart}
\usepackage[margin=19mm]{geometry}
\usepackage{amsmath,amssymb,mathtools}
\usepackage[scr=rsfs,cal=euler]{mathalfa}
\usepackage{xfrac}
\usepackage[unicode,hidelinks,bookmarks=false]{hyperref}
\newcommand{\ie}{i.e.\ }
\newcommand{\cf}{cf.\ }
\newcommand{\resp}{respectively\ }
\newcommand{\Subsection}{\subsection}
\providecommand{\nobreakdash}{\nobreak\hbox{-}}
\setlength{\emergencystretch}{2em}
\multlinegap0pt
\usepackage{bm}
\usepackage{bbm}
\usepackage{dsfont}
\usepackage{xspace}
\usepackage{cancel}
\usepackage{mathtools}
\usepackage{amsthm}
\makeatletter
\@ifundefined{theorem}{\newtheorem{theorem}{Theorem}}{}
\makeatother
\title[Local Linearity of LLMs Enables Activation Steering via Mode]{Local Linearity of LLMs Enables Activation Steering via Model-Based Linear Optimal Control}
\author{Julian Skifstad, Xinyue Annie Yang, Glen Chou}
\date{Source: arXiv:2604.19018v1 (CC BY 4.0)}
\begin{document}
\maketitle

\noindent\emph{Source and licence.} Local Linearity of LLMs Enables Activation Steering via Model-Based Linear Optimal Control. Version arXiv:2604.19018v1 (CC BY 4.0), distributed under the Creative Commons Attribution 4.0 International licence, as recorded in the arXiv licence metadata for this exact version and shown on the abstract page https://arxiv.org/abs/2604.19018v1. Full rights notice: licenses/arxiv-2604.19018-CC-BY-4.0.txt.\par\smallskip

\noindent\textit{Editorial scope.} Counted unit: the Theorem whose source label is \texttt{ap:thm:tracking} in the article (source0.tex lines 919--945, source label \texttt{ap:thm:tracking}), delivered together with the article's own complete original proof of it (source0.tex lines 948--1016, one \texttt{proof} environment of 69 lines). No mathematics is written, repaired or re-derived: statement and proof are the article's own text line for line. Required context retained uncounted: steered\_transformer (lines 229--231); linear\_tf (lines 306--311). Every definition, display and label of those lines is retained unchanged. Omitted from this package and not redistributed: 1--228, 232--305, 312--918, 946--947, 1017--2120. Nothing inside a retained excerpt is omitted or altered: every retained body is delivered exactly as the article's lines are written, so no omission receipt of any kind accompanies this package. Citations: the retained excerpts carry no citation command at all, so no bibliography entry of the article is transcribed and no reference is redistributed. Reference adaptation: none was needed and none was made; every reference inside the delivered unit resolves inside it, so no unresolved reference or citation is produced. Printed slips kept literal: no printed word, symbol, hypothesis or number of the counted statement or of its proof was altered. Layout and class: the article's own class is \texttt{article} with packages including amsmath, microtype, graphicx, subcaption, booktabs, threeparttable, multirow, multicol, caption, placeins, hyperref, float, icml2026, amssymb; the wrapper is the lane's pinned \texttt{amsart} editorial wrapper at 10pt on A4, which restates the theorem-like environments the retained text needs and adds the article's own macro definitions that the retained text needs (none), with extra wrapper packages (bm, bbm, dsfont, xspace, cancel, mathtools). The delivered numbering is the unit's own. Assets: no retained span contains a figure environment, a graphics-inclusion command, a PDF-page inclusion, a table or a TikZ picture; no image file is copied into this package, the package declares no asset dependency and no image file is redistributed with it. Licence: version arXiv:2604.19018v1 of this article is distributed under the Creative Commons Attribution 4.0 International licence, as recorded in the arXiv licence metadata for this exact version and shown on the abstract page https://arxiv.org/abs/2604.19018v1. A full rights notice is delivered at licenses/arxiv-2604.19018-CC-BY-4.0.txt. Characters: every retained excerpt is pure ASCII, so the delivered engine, XeLaTeX with its default Latin Modern text font, reports no missing character.
\begin{equation}\label{steered_transformer}
    z_{k+1} := \rho_k(z_k, u_k) := \phi_k(z_k) + u_k,
\end{equation}

\begin{subequations}\label{linear_tf}
\begin{align}
    \hspace{-5pt}\delta z_{k+1} &\approx A_k \delta z_k + B_k\delta u_k,\\
    A_k &:= (\partial \phi_k/\partial z)\big|_{(\bar z_k=z_{k,+},\bar u_k=0)},\quad B_k := I.\hspace{-5pt}\label{eq:linearized_llm}
\end{align}
\end{subequations}

\begin{theorem}[Closed-Loop Tracking Error Bound]\label{ap:thm:tracking}
Consider the steered LLM dynamics \eqref{steered_transformer}. Assume that $\phi_k$ is twice continuously differentiable for all $k \in \{1,\ldots,\ell\}$. Let $\{(\bar z_k,\bar u_k=0)\}_{k=1}^\ell$ be a nominal trajectory for $\bar z$ constructed via \eqref{linear_tf} and let $\varepsilon_k \in \mathbb{R}^d$ denote the control residual needed to make the mean trajectory $\bar z_k$ satisfy \eqref{steered_transformer}, i.e., $\bar z_{k+1} = \phi_k(\bar z_k) + \varepsilon_k$.
% and define deviations $\delta z_k := z_k - \bar z_k$.
Suppose a linear state-feedback controller $\delta u_k = -K_k \delta z_k$ and define
$\hat A_k := A_k - K_k$, where $A_k := \nabla_z \phi_k(z)|_{z=\bar z_k}$. Assume that for each $k$ there exists a Lipschitz constant $L_k \ge 0$ such that the remainder
\begin{equation}
    r_k(\delta z) := \phi_k(\bar z_k+\delta z) - \phi_k(\bar z_k) - A_k \delta z
\end{equation}
satisfies
$\|r_k(\delta z)\| \le \frac{1}{2} L_k \|\delta z\|^2$ in a neighborhood of $\delta z = 0$, where $\|\cdot\|$ is any vector norm.
Define the closed-loop transition matrices 
\begin{equation}\label{eq:app_transition_matrices}
    \hat\Phi_{k,j} :=
    \begin{cases}
    \hat A_{k-1} \cdots \hat A_j, & k>j, \\
    I, & k=j.
    \end{cases}
\end{equation}

Then for all $k \in \{1,\ldots, \ell\}$,
\begin{subequations}
    \begin{align}
        \Vert \delta z_k\Vert  &\le  \Vert \hat\Phi_{k,1} \Vert \Vert \delta z_1\Vert  + \sum_{i=1}^{k-1} \Vert \hat\Phi_{k,i+1}\Vert \Vert r_i(\delta z_i)+\varepsilon_i\Vert \\
        &\le \|\hat\Phi_{k,1}\|\,\|\delta z_1\| + \sum_{i=1}^{k-1} \|\hat\Phi_{k,i+1}\|\, \Big(\Vert \varepsilon_i\Vert + \frac{L_i}{2}\|\delta z_i\|^2 \Big).
    \end{align}
\end{subequations}
\end{theorem}

\begin{proof}
From the definition of the steered LLM dynamics \eqref{steered_transformer}, the deviation dynamics
satisfy
\begin{align*}
\delta z_{k+1}
&= z_{k+1} - \bar z_{k+1} \\
&= \phi_k(\bar z_k + \delta z_k) + \bar u_k + \delta u_k
   - \bigl(\phi_k(\bar z_k) + \bar u_k + \varepsilon_k\bigr) \\
&= \phi_k(\bar z_k + \delta z_k) - \phi_k(\bar z_k)
   + \delta u_k - \varepsilon_k .
\end{align*}

Using the Taylor expansion of $\phi_k$ about $\bar z_k$,
\[
\phi_k(\bar z_k + \delta z_k)
= \phi_k(\bar z_k) + A_k \delta z_k + r_k(\delta z_k),
\]
which yields
\[
\delta z_{k+1}
= A_k \delta z_k + \delta u_k + r_k(\delta z_k) - \varepsilon_k .
\]

Substituting the feedback law $\delta u_k = -K_k \delta z_k$ gives the exact
closed-loop nonlinear deviation dynamics
\[
\delta z_{k+1}
= \hat A_k \delta z_k + r_k(\delta z_k) - \varepsilon_k .
\]

We now unroll this recursion by induction. For $k=1$,
\[
\delta z_2
= \hat A_1 \delta z_1 + r_1(\delta z_1) - \varepsilon_1
= \hat\Phi_{2,1}\delta z_1
  + \hat\Phi_{2,2} r_1(\delta z_1)
  - \hat\Phi_{2,2} \varepsilon_1 .
\]

Assume the expression holds for some $k$. Then
\begin{align*}
\delta z_{k+1}
&= \hat A_k \delta z_k + r_k(\delta z_k) - \varepsilon_k \\
&= \hat A_k \Bigl(
      \hat\Phi_{k,1} \delta z_1
      + \sum_{i=1}^{k-1} \hat\Phi_{k,i+1} r_i(\delta z_i)
      - \sum_{i=1}^{k-1} \hat\Phi_{k,i+1} \varepsilon_i
    \Bigr)
    + r_k(\delta z_k) - \varepsilon_k \\
&= \hat\Phi_{k+1,1} \delta z_1
  + \sum_{i=1}^{k} \hat\Phi_{k+1,i+1} r_i(\delta z_i)
  - \sum_{i=1}^{k} \hat\Phi_{k+1,i+1} \varepsilon_i ,
\end{align*}
where we used the definitions of the transition matrices
\eqref{eq:app_transition_matrices}. By induction, this expression holds for all
$k = 1,\ldots,\ell$.

Taking norms and applying submultiplicativity yields
\[
\|\delta z_k\|
\le
\|\hat\Phi_{k,1}\|\,\|\delta z_1\|
+ \sum_{i=1}^{k-1} \|\hat\Phi_{k,i+1}\|\,\|r_i(\delta z_i)\|
+ \sum_{i=1}^{k-1} \|\hat\Phi_{k,i+1}\|\,\|\varepsilon_i\| .
\]

Applying the quadratic remainder bound
$\|r_i(\delta z_i)\| \le \tfrac12 L_i \|\delta z_i\|^2$ completes the proof.
\end{proof}

\end{document}
